% CICIL 2026 — Final paper. Compile with:
%   pdflatex cicil_final ; bibtex cicil_final ; pdflatex cicil_final ; pdflatex cicil_final
%
% Content drafted by Romit (2026-07-29) from analysis/paper_bullets.md +
% analysis/human_eval/paper_notes.md. Factual corrections applied during
% LaTeX conversion (each marked with a "FIXED:" comment at the site):
%   - Wixárika has 50 dev images PLUS the 20-image gold-Spanish pilot
%   - retrieval bank sizes updated to the actually-shipped banks
%     (Guaraní 15,494 — MultiScript30k excluded on license grounds; Bribri 8,297)
%   - caption-length column recomputed on dev (Nahuatl 8.3, Wixárika 17.0)
%   - abstract culture-name range 18–37 (student) / 29–38 (teacher)
% TODO-TONIGHT markers = numbers that land after the 2026-07-29 overnight runs.
\pdfoutput=1
\documentclass[11pt]{article}
\usepackage{ACL2023}
\usepackage{times}
\usepackage{latexsym}
\usepackage[T1]{fontenc}
\usepackage[utf8]{inputenc}
\usepackage{microtype}
\usepackage{inconsolata}
\usepackage{booktabs}
\usepackage{graphicx}

\title{Culturally-Aware Image Captioning for Indigenous Languages of the Americas:\\
A Retrieval-Grounded and Verification-Oriented Vision-to-Translation Pipeline}

\author{Romit Basak \and Mehak Santosh Talmale \and Nandita Mahendra \and Tisha Patel \\
  Khoury College of Computer Sciences, Northeastern University \\
  \texttt{\{basak.r, talmale.m, mahendra.na, patel.tisha\}@northeastern.edu}}

\begin{document}
\maketitle

\begin{abstract}
Generic vision--language models (VLMs) frequently fail to preserve culturally
specific information present in images, creating a bottleneck for multilingual
image captioning systems targeting Indigenous languages. In the AmericasNLP
2026 Cultural Image Captioning for Indigenous Languages (CICIL) shared task,
captions are generated directly in Indigenous languages of the Americas from
culturally situated images. Existing systems typically generate an
intermediate Spanish caption with a generic VLM and then translate it using
retrieval-augmented language models, but culturally salient visual concepts
are often lost before translation begins. We investigate whether
retrieval-grounded cultural reasoning can improve this stage.

We present a two-stage captioning pipeline that augments vision processing
with cultural retrieval and verification. Stage~1 uses structured cultural
visual question answering (VQA), Wikipedia-derived retrieval, and a
retrieval-grounded interrogation loop inspired by caption-questioning
frameworks. Retrieved cultural hypotheses are verified through closed visual
questions before being incorporated into Spanish descriptions. Stage~2
performs culturally indexed retrieval over parallel corpora and translates
the Spanish intermediate into the target language using Gemini~2.5~Flash.

Across five Indigenous languages (Guaraní, Bribri, Yucatec Maya, Wixárika,
and Nahuatl), retrieval substantially increases explicit cultural references,
% FIXED: was "18–39"; student arm range is 18–37/50, teacher 29–38/50
raising culture-name rates from near-zero to 18--37 of 50 dev captions
(29--38 for a 7B teacher). However, corresponding chrF++ improvements are
small and inconsistent. Through human evaluation, retrieval audits, decoding
ablations, and extensive qualitative analysis, we show that standard
reference-based metrics are largely insensitive to cultural correctness,
hallucinations, retrieval-induced errors, and faithfulness improvements. We
further demonstrate that verification ability is capability-gated: a 2B model
fails as a verifier even under closed questioning, while a local 7B model
reliably distinguishes true from false claims. Our results suggest that
culturally grounded captioning is constrained less by translation quality
than by cultural retrieval, verification capacity, and evaluation
methodology.
\end{abstract}

\section{Introduction}
The AmericasNLP 2026 Cultural Image Captioning for Indigenous Languages
(CICIL) shared task \citep{bui2026cicil} asks systems to generate image
captions directly in Indigenous languages of the Americas. Unlike
conventional image captioning benchmarks, CICIL emphasizes cultural
grounding: images depict artifacts, ceremonies, landscapes, and community
practices that are often meaningful only within specific cultural contexts.

The strongest systems in recent shared-task settings adopt a two-stage
pipeline \citep{dhawan2026gators}. A vision--language model first produces an
intermediate Spanish description, which is then translated through
retrieval-augmented language models into the target language. While
effective, this design inherits a fundamental limitation: any cultural
information omitted in the Spanish intermediate cannot be recovered
downstream.

Our work starts from a simple observation. Culture recognition from pixels
alone is an ill-posed inverse problem. Many visual motifs recur across
unrelated traditions, geographic regions, and historical periods. Ñandutí
lace patterns, for example, descend from broader Tenerife lace traditions and
cannot always be uniquely identified from image evidence alone. Cultural
attribution therefore requires not only visual evidence but also contextual
priors. The CICIL dataset already provides one such prior through image-level
cultural metadata, yet most existing pipelines do not incorporate this
information into visual reasoning.

We investigate whether retrieval-grounded cultural reasoning can improve
culturally situated caption generation. Our initial experiments with cultural
VQA prompting showed only marginal gains, motivating a stronger
retrieval-based approach. We therefore introduce a retrieval-grounded
captioning architecture in which external cultural knowledge is supplied at
inference time and verified through visual questioning before entering the
final caption.

The paper addresses three research questions:

\paragraph{RQ1:} Does culturally grounded Stage-1 processing improve caption
quality compared with generic vision--language baselines?
\paragraph{RQ2:} How does performance vary across languages with different
resource availability?
\paragraph{RQ3:} Which categories of cultural information are most difficult
to ground reliably?

Our contributions are:
\begin{itemize}
  \item A retrieval-grounded cultural captioning pipeline combining visual
    reasoning, cultural retrieval, and verification.
  \item Evidence that retrieval, rather than model scale alone, is the
    primary constraint on cultural naming.
  \item A capability analysis showing that verification emerges only above a
    minimum model scale.
  \item An empirical demonstration that chrF++ fails to reflect many
    important dimensions of cultural caption quality.
  \item Extensive qualitative analysis of retrieval failures, hallucinations,
    calibration behavior, and evaluation limitations.
\end{itemize}

\section{Related Work}

\subsection{Cultural Grounding in Vision--Language Models}
Vision--language models inherit cultural biases from predominantly Western
web-scale training data. CLIP-style contrastive pretraining
\citep{radford2021clip} produces powerful multimodal representations but
remains sensitive to geographic and cultural imbalances. Benchmarks such as
MaRVL \citep{liu2021marvl} demonstrate large performance drops when visual
reasoning is evaluated across cultures rather than within familiar
English-centric settings.

More recently, CulturalVQA \citep{nayak2024culturalvqa} and related
benchmarks \citep{romero2024cvqa} have shown that state-of-the-art systems
still struggle with culturally specific visual concepts. Fine-tuning
approaches such as CultureCLIP \citep{huang2025cultureclip} improve cultural
recognition through synthetic contrastive examples, but such methods require
large curated image collections. Under CICIL's small-data constraints,
prompting and retrieval become more practical alternatives.

The CIC framework \citep{yun2024cic} introduced a VQA-before-captioning
paradigm in which a model first identifies cultural elements and then
synthesizes a caption. Our work adopts this idea but combines it with
retrieval-grounded reasoning and downstream translation.

\subsection{Cultural Retrieval}
Retrieval-augmented systems have become increasingly important for
low-resource NLP \citep{lewis2020rag}. Instead of requiring all knowledge to
be encoded in model parameters, retrieval supplies relevant information
dynamically at inference time.

For cultural grounding, this distinction is particularly important. In our
experiments, even a 7B teacher model rarely names cultures unaided, while
retrieval dramatically increases culture naming rates. This suggests that
cultural knowledge is not merely missing from smaller models but is
insufficiently accessible even in larger ones.

A key insight from our findings is that retrieval quality and retrieval usage
are distinct. We observed examples where retrieval correctly retrieved
relevant cultural material but the model copied irrelevant metadata instead
of grounding itself in the image.

\subsection{Interrogation and Verification}
Our retrieval-grounded interrogation loop follows the ChatCaptioner family of
approaches \citep{zhu2023chatcaptioner}, where an LLM asks questions and a
vision model answers them before caption synthesis. Related work includes
Video ChatCaptioner \citep{chen2023videochatcaptioner}, IdealGPT
\citep{you2023idealgpt}, and visually grounded dialogue systems such as
Visual Dialog \citep{das2017visualdialog} and GuessWhat?!
\citep{devries2017guesswhat}.

The verification component is closely related to Chain-of-Verification
\citep{dhuliawala2024cove} and Woodpecker \citep{yin2023woodpecker}. However,
unlike prior work, we explicitly measure \emph{when} verification works.

A central finding of this paper is that verification itself is
capability-gated. A 2B model answered ``yes'' to most verification prompts
regardless of correctness, while a local 7B model reliably rejected false
claims. Verification therefore appears to require a minimum level of visual
reasoning capacity before verification-based pipelines become effective.
Notably, ChatCaptioner's own error analysis already attributes 94\% of its
incorrect captions to wrong answers from the answerer
\citep{zhu2023chatcaptioner}; our results characterize where that answerer
capacity turns on.

\section{Data}
% FIXED: Wixárika has a 50-image dev set like the other languages; the
% 20-image pilot (with gold Spanish) is an ADDITIONAL split.
The CICIL development data are extremely small: each of the five languages
provides 50 development images. Wixárika additionally provides a 20-image
pilot set that is uniquely valuable because it includes gold Spanish
captions (mean 13.7 words), enabling direct evaluation of Stage~1.

\begin{table}[t]
\centering
\small
\begin{tabular}{lrrr}
\toprule
Language & Dev & Caption len. & Bank pairs \\
\midrule
Guaraní       & 50 & 21.6 & 15{,}494 \\
Bribri        & 50 & 14.9 & 8{,}297 \\
Yucatec Maya  & 50 & 11.3 & 14{,}332 \\
Nahuatl       & 50 & 8.3  & 16{,}119 \\
Wixárika      & 50$^{\dagger}$ & 17.0 & 9{,}940 \\
\bottomrule
\end{tabular}
\caption{Development data and Stage-2 retrieval banks. ``Caption len.'' =
mean gold target-caption length (words). $^{\dagger}$Plus a 20-image pilot
set with gold Spanish captions. Banks are built from AmericasNLP shared-task
corpora \citep{mager2021americasnlp,ebrahimi2023americasnlp}, the
Wixárika--Spanish corpus \citep{mager2018wixarika}, and the YUA-ES
Communicative Contexts Corpus \citep{molina2026yuaes}.}
% FIXED: bank sizes were stale in the draft (Guaraní 53,183 → 15,494:
% MultiScript30k is excluded on license grounds; Bribri 7,506 → 8,297).
\label{tab:data}
\end{table}

Resource availability varies dramatically across languages
(Table~\ref{tab:data}), creating a natural testbed for understanding how
retrieval affects low-resource systems.

\section{Method}

\subsection{Stage 1: Cultural VQA}
The original Stage~1 system decomposes image understanding into four cultural
categories: (i)~ceremony and ritual, (ii)~material culture, (iii)~landscape
and geography, and (iv)~kinship and community. A small VLM
(SmolVLM-2B; \citealp{marafioti2025smolvlm}) answers category-specific
questions and then synthesizes a Spanish description from the answers,
following the VQA-before-captioning paradigm \citep{yun2024cic}.

Our initial experiments revealed that verbose cultural explanations lower
chrF++ because references are short. A constrained single-sentence synthesis
restored performance while preserving cultural reasoning.

\subsection{Retrieval-Grounded Cultural Context}
We augment Stage~1 with two cultural retrieval channels.

\paragraph{Image retrieval.} We construct SigLIP \citep{zhai2023siglip}
embedding indices over culture-specific Wikimedia Commons image collections.
Retrieved neighbors provide visual examples associated with the target
culture.

\paragraph{Text retrieval.} We build culture-specific Wikipedia knowledge
banks and retrieve relevant passages using multilingual MiniLM embeddings
\citep{reimers2020multilingual}. These passages serve as cultural context
during synthesis.

Retrieved evidence is associated with confidence bands: \emph{strong match}
($\geq 0.80$ similarity) and \emph{possible match} ($0.55$--$0.80$). Only
strong evidence supports direct assertions; possible matches are hedged
explicitly.

\subsection{Multi-Agent Interrogation}
The final version of Stage~1 adopts a retrieval-grounded interrogation
architecture:
\begin{enumerate}
  \item Generate a base caption.
  \item Retrieve candidate cultural concepts.
  \item Generate closed verification questions.
  \item Answer questions using a VLM.
  \item Assemble a final caption from verified evidence.
\end{enumerate}
The questioner and assembler are text-only LLMs; the answerer is a local 7B
VLM (Qwen2.5-VL; \citealp{qwen2025qwen25vl}). This design transforms open
cultural speculation into a sequence of verifiable claims.

\subsection{Capability-Gated Verification}
Verification performance differed dramatically by model size.

The 2B model displayed extreme yes-bias, producing 82 affirmative responses
out of 82 verification attempts. On planted decoys, it accepted clearly false
claims and sometimes contradicted its own rationale.

In contrast, a local 7B model correctly rejected all negative controls and
successfully verified true positives.

These results indicate that visual verification appears only above a
capability threshold.

\subsection{Stage 2 Translation}
The final Spanish description is translated using retrieval-augmented
Gemini~2.5~Flash \citep{comanici2025gemini}.

We index parallel corpora using multilingual MiniLM embeddings and FAISS
\citep{johnson2019faiss} and retrieve top-$k$ Spanish--Indigenous examples.
Retrieved pairs are supplied as demonstrations during translation.

Temperature 0.7 decoding was adopted after ablation experiments revealed
severe degeneration under greedy decoding.

\section{Results}

\subsection{Stage 1 Evaluation}
Using the Wixárika pilot's gold Spanish captions, we score Stage-1 Spanish
intermediates with chrF++ \citep{popovic2017chrfpp}:

\begin{table}[h]
\centering
\small
\begin{tabular}{lr}
\toprule
Configuration & chrF++ \\
\midrule
Generic                 & 21.0 \\
Cultural VQA (verbose)  & 16.2 \\
Cultural VQA (concise)  & 20.9 \\
\bottomrule
\end{tabular}
\caption{Spanish-side proxy evaluation (Wixárika pilot, $n{=}20$).}
\label{tab:proxy}
\end{table}

Concise cultural prompting reaches parity with generic captioning while
reducing internal inconsistencies.

\subsection{Human Evaluation}
Single-annotator results over 15 images (0--2 rubric per dimension), judged
on English pivot translations of the Spanish captions
\citep{hendy2023gpt,kocmi2023wmt}:

\begin{table}[h]
\centering
\small
\begin{tabular}{lrr}
\toprule
Dimension & Generic & Cultural \\
\midrule
Cultural accuracy & 0.00 & 0.13 \\
Faithfulness      & 1.73 & 1.53 \\
Fluency           & 1.93 & 1.87 \\
\bottomrule
\end{tabular}
\caption{Human evaluation of Stage-1 arms. Preference outcomes: tie 60\%,
cultural preferred 27\%, generic preferred 13\%.}
\label{tab:human}
\end{table}

% TODO-TONIGHT: round-4 human eval (40 images, 5 cultures, two annotators,
% baseline vs agent-local v4.5 finals) + inter-annotator kappa land tomorrow
% morning — add a second human-eval table here.

The main conclusion is that cultural prompting alone generates slightly more
cultural content but remains largely indistinguishable from generic
captioning.

\subsection{Decoding Ablation}
Greedy decoding created severe repetition loops:

\begin{table}[h]
\centering
\small
\begin{tabular}{lrr}
\toprule
Language & Degen.\ before & Degen.\ after \\
\midrule
Bribri   & 64\% & 28\% \\
Wixárika & 86\% & 30\% \\
\bottomrule
\end{tabular}
\caption{Degenerate-caption rate under greedy decoding (before) vs
temperature 0.7 with fixed seed (after), 50 dev images per language.}
\label{tab:decoding}
\end{table}

Switching to temperature-based sampling reduced degeneration dramatically and
improved chrF++ without harming stronger languages.

\subsection{Retrieval Results}
Retrieval dramatically increases explicit cultural references
(Table~\ref{tab:rates}).

\begin{table}[t]
\centering
\small
\begin{tabular}{lrrr}
\toprule
\multicolumn{4}{c}{\textbf{Culture-name rate} (of 50)} \\
Language & No RAG & RAG student & RAG teacher \\
\midrule
Guaraní   & 3 & 37 & 37 \\
Wixárika  & 0 & 21 & 38 \\
Maya      & 1 & 32 & 38 \\
Nahuatl   & 1 & 18 & 29 \\
Bribri    & 0 & 34 & 37 \\
\midrule
\multicolumn{4}{c}{\textbf{Concept rate} (of 50)} \\
Language & No RAG & RAG student & \\
\midrule
Wixárika  & 0 & 24 & \\
Bribri    & 5 & 25 & \\
Guaraní   & 1 & 5  & \\
Nahuatl   & 0 & 1  & \\
Maya      & 0 & 0  & \\
\bottomrule
\end{tabular}
\caption{Cultural specificity with and without retrieval. ``Culture-name
rate'' counts captions naming the target culture; ``concept rate'' counts
captions naming a specific cultural concept (artifact, dance, site).}
\label{tab:rates}
\end{table}

Retrieval therefore succeeds at increasing cultural specificity even when
chrF++ remains nearly unchanged.

\subsection{Capability Ladder}
Wixárika pilot results across system versions:

\begin{table}[h]
\centering
\small
\begin{tabular}{lr}
\toprule
System & chrF++ \\
\midrule
SmolVLM-RAG        & 18.44 \\
RAG distill        & 18.79 \\
Verify-RAG$^{*}$   & 20.50 \\
Local 7B base      & 21.47--21.75 \\
Agent v4.2         & 21.73 \\
Agent v4.3 (local) & 21.70 \\
% TODO-TONIGHT: add final agent row (v4.3.1 transcripts + v4.5 whitelist
% aggregation, fully local) when tonight's runs land.
Gemini direct      & 22.71 \\
\bottomrule
\end{tabular}
\caption{Capability ladder on the Wixárika pilot (chrF++ vs gold Spanish).
$^{*}$The Verify-RAG score is inflated for the wrong reason: its 2B verifier
accepted every claim (\S4.4), so its gain reflects longer, more assertive
captions, not verification.}
\label{tab:ladder}
\end{table}

The remaining gap to frontier performance is roughly 1--1.5 chrF++.

\section{Analysis}

\paragraph{RQ1: Does cultural grounding help?}
Cultural prompting alone does not improve end-to-end performance. Retrieval
substantially increases cultural specificity but produces limited metric
gains. Human evaluation suggests that cultural prompting yields marginally
richer cultural descriptions at the cost of slight reductions in
faithfulness.

\paragraph{RQ2: Resource effects.}
Resource disparities appear throughout the pipeline. Bribri has the smallest
retrieval resources and remains difficult across all systems. However,
retrieval produced some of the largest relative gains in cultural concept
rates, suggesting that low-resource settings may benefit disproportionately
when even modest external knowledge becomes available.

\paragraph{RQ3: Which cultural concepts ground reliably?}
Artifacts and distinctive cultural objects appear most retrievable. Landscape
concepts proved much harder: several initially promising examples collapsed
under manual audit because retrieval supplied plausible but unsupported
regional labels. This finding motivated stricter auditing procedures and
highlights the danger of equating concept naming with genuine grounding.

\section{Conclusion}
We presented a retrieval-grounded cultural image captioning pipeline for
Indigenous languages of the Americas. Cultural prompting alone provided
minimal gains, but retrieval substantially improved cultural specificity. At
the same time, our experiments reveal that cultural captioning is ultimately
constrained not by translation but by retrieval quality, verification
capability, and evaluation methodology.

Most importantly, we show that standard reference-based metrics fail to
capture many of the qualities that matter in culturally situated captioning
systems. Future work should therefore prioritize verified cultural
grounding, community-informed evaluation, and retrieval methods that provide
trustworthy contextual knowledge rather than encouraging plausible cultural
guessing.

\section*{Limitations}
This work has several limitations.

First, reference-based metrics are largely insensitive to cultural
correctness. Multiple faithfulness improvements, hallucination fixes, and
verification gains produced almost no change in chrF++. Conversely,
degenerate repetition loops still received nontrivial scores.

Second, cultural attribution is inherently underdetermined from image
evidence alone. Many images require contextual priors that are unavailable
from pixels.

Third, retrieval quality remains an open problem. Manual inspection revealed
that many medium-confidence retrievals matched scene categories rather than
specific cultural artifacts. Concept-rate statistics may therefore reflect a
mixture of genuine grounding and retrieval-induced label copying.

Fourth, our human evaluation was conducted by researchers who were neither
community members nor speakers of the target languages. Cultural judgments
should therefore be interpreted as measurements of specificity rather than
definitive correctness.

Fifth, gold captions frequently contain nonvisual contextual knowledge.
Approximately two-thirds of the Wixárika pilot references include information
such as event purpose, seasonal context, or community-specific background
that cannot be inferred directly from pixels. This imposes a structural
ceiling on image-only captioning systems.

Finally, we did not conduct target-language human evaluation due to the
absence of available native-speaker annotators during the project period.

\section*{Ethics Statement}
Our goal is to support Indigenous-language technology without replacing
community expertise. Generated captions are not authoritative cultural
interpretations and should not be deployed in community-facing contexts
without review by native and heritage speakers.

We adopt a conservative policy that prohibits naming photographed individuals
unless their identity is explicitly present in visible image text. This
decision follows observed cases where retrieval systems incorrectly attached
real cultural figures to unrelated individuals.

All datasets and retrieval sources were used according to their published
licenses. Image retrieval collections were harvested only from appropriately
licensed Wikimedia Commons content, and translation was performed through a
paid-tier API configuration compatible with CICIL licensing requirements.
The Maya retrieval bank uses the YUA-ES Communicative Contexts Corpus
\citep{molina2026yuaes}, cited per its release requirements.

\bibliography{custom}
\bibliographystyle{acl_natbib}

\end{document}
